% =============================================================================
\section{Conclusions and Future Work}
\label{sec:conclusions}
% =============================================================================

\subsection{The Case for a Hybrid Pipeline}

The result we most want to draw out is the shape of the system rather than any
single stage. Into Frame mixes learned models and classical graphics roughly
two-to-one by stage count, and neither part would produce this output alone.

Generative models make the problem tractable at all. Five sixths of the sphere and
the ground behind every object are never observed, and no classical reasoning
recovers content absent from the input. Outpainting, foreground removal,
single-view reconstruction and video generation supply it.

What they do not supply is structure. A generative system asked for a scene returns
something viewable: a panorama, a mesh, a radiance field. It does not return a
collidable surface, a population whose density can be sampled, a water plane
animated separately from its bed, or a mesh with a skeleton in it. Every property
that makes the output walkable rather than watchable comes from the classical side:
a harmonic solve that provably never modifies measured ground, fluvial incision,
point-process synthesis fitted to observed statistics, pattern texturing, and a
three-bone rig driven by a shared wind field. Of the 38 active stages, 12 use no
learned model at all, and those are disproportionately the ones producing physical
behavior.

The division of labor is the point. Learned components answer \emph{what is there},
including where nothing was observed; classical components decide \emph{how it
behaves}, under constraints a generative model cannot honor: measured terrain must
survive completion unchanged, water must sit below its bed, a meadow of thousands
of plants must remain one instanced draw. Treating those as requirements rather
than as things to be learned is what turns a picture of a place into a place.

It also explains our failures. As \cref{sec:results} sets out, the surviving defects
are mostly not model failures but assumption failures, where a step correct for one
scene category is applied silently to another. That is the more tractable problem:
an assumption can be identified, measured, and gated.

\subsection{Future Work}

\paragraph{Explicit scene classification.} A first-class scene-type estimate
(natural landform, urban, coastal, interior) consumed by every stage that
currently embeds an implicit assumption would address the largest class of failure
we observed. The system is biased toward outdoor natural scenes, and several
assumptions do not hold outside them: that the sky boundary is the crest of ground,
that a large uniform region is terrain, that vegetation is the dominant content.

\paragraph{Confidence propagation.} Stages produce point estimates. Propagating
uncertainty would let downstream stages weigh evidence rather than accept it, and
would let the system report where it is guessing.

\paragraph{Better single-view reconstruction for thin structures.} Vegetation is
the dominant content of natural scenes and the weakest part of our object track;
species-aware priors~\cite{lee2024tree} are the obvious direction.

\paragraph{Joint refinement.} The pipeline is strictly feed-forward. A closing
optimization that adjusted depth, geometry and placement jointly against
photometric consistency could correct compounded error that no single stage can
see.

\paragraph{Perceptual evaluation.} Our evaluation is diagnostic: it measures where
the system is wrong, not how the result feels to be inside. A user study of
presence, scale perception and navigation comfort would test what the system is
actually for, and is the clearest gap in our rigor.

\paragraph{Audio and interaction.} The scene is something to move through rather
than to act on. Spatialized audio generated jointly with the
scene~\cite{jin2026sonoworld}, and ways for the user to disturb the environment
rather than only observe it, would extend interactivity in the direction this
framing implies.

\subsection{Closing Remarks}

Into Frame demonstrates that a single photograph can become an interactive
environment, with terrain, objects, ground cover, water and lighting reconstructed
as distinct physical components rather than as one representation serving every
purpose. That decomposition buys the depth of the result: each component is
reconstructed by a method suited to its role and then behaves according to it.

The remaining failures are mostly not failures of any model, but assumptions that
held for outdoor natural landscapes and did not extrapolate beyond them. That is an
encouraging place to stop: the boundaries are legible and measurable, and addressing
them means making implicit commitments explicit rather than waiting for better
models.
